\section{Special ase: Smooth DCA Lower Bound}
\label{Appendix: lower bound of DCA}
The preceding upper bound has an unavoidable $\delta_g^{-1}$ dependence in worst-case first-order complexity. The cleanest way to see this is a smooth embedding reduction.

\begin{lemma}[Approximate Fenchel criticality implies smooth stationarity]
Let $h$ be convex and $L$-smooth. If $v\in\partial_\delta h(x)$, then
\begin{equation}
    \|v-\nabla h(x)\|\le \sqrt{2L\delta}.
    \label{eq:fenchel-gap-smooth-conversion}
\end{equation}
Consequently, if $\phi$ and $g$ are convex $L$-smooth and
\begin{equation}
    a\in\partial_{\delta_\phi}\phi(x),
    \quad
    b\in\partial_{\delta_g}g(x),
    \quad
    \|a-b\|\le\eta,
\end{equation}
then
\begin{equation}
    \|\nabla(\phi-g)(x)\|
    \le
    \eta+
    \sqrt{2L\delta_\phi}+
    \sqrt{2L\delta_g}.
    \label{eq:relaxed-to-smooth-stationarity}
\end{equation}
\end{lemma}

\begin{proof}
For convex $L$-smooth $h$, the conjugate $h^*$ is $1/L$-strongly convex on its effective domain. Since
\begin{equation}
    \Phi_h(x,v)=h(x)+h^*(v)-\langle v,x\rangle
    =D_{h^*}(v,\nabla h(x)),
\end{equation}
we have
\begin{equation}
    \Phi_h(x,v)\ge \frac{1}{2L}\|v-\nabla h(x)\|^2.
\end{equation}
The implication \eqref{eq:fenchel-gap-smooth-conversion} follows from $\Phi_h(x,v)\le\delta$. Applying this to $\phi$ and $g$ and using the triangle inequality proves \eqref{eq:relaxed-to-smooth-stationarity}.
\end{proof}

\begin{theorem}[Smooth embedding lower bound]
Fix a component oracle that returns
\begin{equation}
    (\phi(x),g(x),\nabla\phi(x),\nabla g(x)).
\end{equation}
There exists a smooth DC subclass with both components convex and smooth such that any algorithm that guarantees a point satisfying
\begin{equation}
    \delta_\phi,\delta_g\le c\eps^2/L,
    \qquad
    \eta\le c\eps
\end{equation}
for a sufficiently small numerical constant $c>0$ requires
\begin{equation}
    \Omega(\Delta L\eps^{-2})
    =\Omega(\Delta/\delta)
\end{equation}
component-oracle calls, where $\delta\asymp\eps^2/L$.
\end{theorem}

\begin{proof}
Let $\widehat F$ be a smooth nonconvex hard instance with $L_0$-Lipschitz gradient and initial gap at most $\Delta$, for which any first-order method needs $\Omega(\Delta L_0\eps^{-2})$ queries to find a point satisfying $\|\nabla\widehat F(x)\|\le\eps$; see \citet{carmon2020lower1,carmon2020lower2}. Define
\begin{equation}
    q(x):=\frac{L_0}{2}\|x\|^2,
    \qquad
    \phi(x):=\widehat F(x)+q(x),
    \qquad
    g(x):=q(x).
\end{equation}
Since $\widehat F$ is $L_0$-smooth, $\widehat F+q$ is convex. Moreover, $\phi$ is $2L_0$-smooth and $g$ is $L_0$-smooth. Set the component smoothness budget $L=2L_0$. The difference satisfies $\phi-g=\widehat F$.

The component oracle is equivalent to the standard first-order oracle for $\widehat F$ because
\begin{equation}
    \widehat F(x)=\phi(x)-g(x),
    \qquad
    \nabla\widehat F(x)=\nabla\phi(x)-\nabla g(x),
\end{equation}
while the reverse simulation is obtained by adding the known quadratic $q$. Finally, Lemma~1 implies that any returned point with $\delta_\phi,\delta_g\le c\eps^2/L$ and $\eta\le c\eps$ is an $\eps$-stationary point of $\widehat F$, after decreasing $c$ if needed. A faster algorithm would contradict the smooth lower bound.
\end{proof}

\begin{remark}[Condition-number phrasing]
If the paper states the relaxed-to-smooth conversion through a strong convexity constant $\gamma$ of $g$ and a smoothness constant $L$, a conservative regime is
\begin{equation}
    \delta_g\lesssim \gamma\eps^2/L^2.
\end{equation}
Under that normalization, the same smooth embedding gives a lower bound often written as
\begin{equation}
    \Omega\left(\frac{\gamma}{L}\cdot\frac{\Delta}{\delta_g}\right),
\end{equation}
which matches the $O(\Delta/\delta_g)$ upper bound up to the condition-number factor $L/\gamma$. If the smooth DC subclass allows $g$ to be a quadratic with curvature of order $L$, the clean $\Omega(\Delta/\delta)$ form applies.
\end{remark}